% !TEX root = ../main.tex
\chapter{Conclusiones}\label{cap:conclusiones}

El trabajo partió de una pregunta: si los algoritmos cuánticos de la era NISQ ofrecen una ventaja medible frente a los clásicos en problemas de ciencia de datos. Para responderla se compararon seis modelos de clasificación sobre Iris, Bars and Stripes e Ising y seis solvers de optimización sobre TSP, MDS y MaxCut, con un protocolo común y la representación como eje de lectura. Las conclusiones siguen de los resultados del Capítulo~\ref{cap:evaluacion}.

\section{Conclusiones}

\begin{enumerate}
\item \textbf{En clasificación no se observa ventaja cuántica.} El SVM con kernel RBF obtiene el mejor F1 o empata con el mejor en los tres datasets (0.967 en Iris, 0.870 en Bars and Stripes, 0.985 en Ising) y es el más rápido de entrenar. Ningún modelo cuántico lo supera. La media cuántica de accuracy queda 0.083 y 0.101 por debajo de la clásica en Iris y Bars and Stripes, y la iguala en Ising.

\item \textbf{El QSVM con mapa ZZ solo alcanza al clásico en el dataset más fácil.} Empata con el SVM en Ising (F1 0.985), donde los seis modelos alcanzan al menos 0.95 de accuracy, y cae a 0.730 en Iris y 0.758 en Bars and Stripes, con el mayor gap de generalización del corpus en Iris (0.233). El experimento no aísla la causa del empate en Ising: el mapa actúa sobre componentes principales de las amplitudes, no sobre los espines de la cadena.

\item \textbf{Los modelos variacionales se acercan al clásico sin superarlo.} La QNN híbrida es el modelo cuántico más regular (F1 0.933, 0.815 y 0.937, con gap de generalización $\le 0.022$). El VQC iguala a los árboles en Iris (0.950), queda el último en Bars and Stripes (0.688) y es el más caro de entrenar en los tres datasets.

\item \textbf{En optimización, QAOA y VQE reproducen el óptimo, pero no mejoran a los métodos clásicos.} Los métodos exactos alcanzan $C^\star$ en las tres instancias, la enumeración en menos de un milisegundo. El annealing simulado lo hace sobre la misma formulación QUBO en milisegundos (éxito 0.9 en TSP y 1 en MDS y MaxCut). QAOA y VQE llegan al óptimo en TSP y MaxCut, y VQE también en MDS, con entre 0.6 y 2~s de simulación. En MDS una de las ocho semillas de QAOA devuelve una solución infactible (éxito 0.875, 12.6~s): termina con mayor energía esperada y probabilidad nula de medir el óptimo, mientras que el annealing sobre el mismo QUBO acierta siempre. Como en TSP y MDS ambos leen el mejor de miles de bitstrings, su éxito en esos casos no prueba que el circuito concentre probabilidad en el óptimo.

\item \textbf{El régimen estudiado mide proximidad al clásico, no ventaja.} Con 4 qubits en clasificación, como mucho 6 en optimización y simulación de estado exacta, los resultados no permiten afirmar ni descartar una ventaja a escala \citep{preskill2018, thanasilp2022}. Sí permiten afirmar que, en este rango, ningún modelo o solver cuántico supera al mejor clásico, y que la paridad de calidad, cuando aparece, va acompañada de un coste temporal mayor.
\end{enumerate}
